While our findings corroborate prior work and provide strong evidence
for the proposed scaling laws in $\tepoch$, $\bopt$, and $\bcrit$,
there are several limitations that merit further study.

\paragraph{EMA perspective}

As the EMA perspective regards parameters $y_t$ as a function of
updates $x_t$, it fails to account for $x_t$ actually depending on
earlier values of $y_t$ (e.g., $y_{t-1}$). Yet although this
perspective has formal limitations, we nevertheless find it a
useful \emph{conceptual} model of training, as it predicts behavior
that is supported by experiments.

\paragraph{Optimization and training setup}

Our work focuses on AdamW (the standard optimizer for LLM training).
While the EMA perspective applies directly to other optimizers that
use decoupled weight decay, such as Sophia~\citep{liu2023sophia} and
MuonClip~\citep{kimi2025k2}, it may not apply to approximate second
order methods, e.g., Shampoo~\citep{gupta2018shampoo}.  However, it
can be used when applying AdamW (and related optimizers) in Shampoo's
eigenbasis, which was shown to be effective in
SOAP~\citep{vyas2024soap}.

We present results with a single (standard) learning rate schedule.
Our method for obtaining $\cbs$ estimates would be quite efficient
with a warmup-stable-decay (WSD)
schedule~\citep{hu2024minicpm,wen2024understanding}, as we could
perform a single training run with each batch size, but decay at
various milestones in order to get points along the scaling law,
essentially following the approach in~\citet{hagele2024scaling}, but
with separate laws for each batch size.

We used the maximal update parameterization in all experiments, which
generates a learning rate $\eta$ adjustment for each model width.  Our
results suggest this approach enables good models at arbitrary $N$,
$D$, and $B$ when combined with adjustments to $\lambda$.  This
strategy is informed by our experiments comparing re-adjusting $\eta$
vs.\ $\lambda$ in \cref{sec:tema_results}. However, it is not
feasible, at this scale, to verify whether substantially better models
could be obtained by sweeping the full cross-product of $\eta$ and
$\lambda$ values.

Our study specifically focuses on the practically important setting of
single-epoch LLM pre-training.  \citet{wang2024how} indeed noted
differences in optimal $\tema$ when using multi-epoch training,
possibly due to data repetition. Reconciling these differences by
isolating the effects of repetition versus scale is an interesting
follow-up direction.

Here we only experimented with a single dataset, vocabulary, and
context length.  We obtained a similar $\cbs$ scaling law
to \citet{zhang2024how}, but it would be interesting to see if
differences in the coefficient of our power laws could be attributed
to specific differences in approach (e.g., differences in dataset,
context length, learning rate schedule, use of weight decay,
etc.).  \cref{sec:appendix_zhang} has further discussion of
differences with \citet{zhang2024how}.

We have also not explored how changes in numerical precision could
affect scaling laws.  Recent work~\citep{kumar2024scaling} showed
that, in terms of scaling laws, lower precision reduces the model's
\emph{effective} parameter count.  This suggests precision would
have no impact on scaling of $\bopt$ or $\bcrit$, which do not scale
in $N$.  Lower precision, however, could increase the \emph{effective}
TPP (via smaller effective $N$), thereby altering $\temaopt$.

\paragraph{Small batches, large batches, and dynamic batch sizing}

We consistently find that smaller and smaller batches do not grow
asymptotically closer to $\dmin$, as predicted by theory, but
eventually degrade in loss.  One possibility is that $\lambda$ tuning
is not sufficient with very small $B$, and further tuning of other
hyperparameters may be needed, such as the Adam $\beta$ parameters (as
suggested in recent
work~\citep{porian2024resolving,zhang2024how,marek2025small}).
%
Some preliminary tests using the $\beta_2$ scaling rule
from \citet{marek2025small} showed loss improvements at small $B$.
%
Since we are unlikely to train with small batches at scale, and using
them even with smaller LLMs significantly impairs our ability to train
both efficiently and quickly, it is unfortunately difficult to justify
further exploration in this direction.

Regarding large batches, our methods do not account for the many
practical systems-related issues, including bandwidth and
communication overheads, memory limits of hardware, synchronization
delays, etc.  
%
Moreover, as batch sizes increase, techniques such as optimizer
sharding may be needed, which further complicate performance
model~\citep{almazrouei2023falcon}.
%
Our scaling laws do, however, explicitly define a practically relevant
regime of training batch sizes: $\bopt \le
B \le \bcrit$. Practitioners can leverage this identified regime
alongside system-specific profiling (e.g., evaluating utilization at
various batch sizes) to select optimal settings balancing algorithmic
and systems constraints.

Exploring optimal dynamic batch sizing is a natural future direction
for our work.  While the potential gains were found to be small in
theory by \citet{mccandlish2018empirical}, more recent work has found
significant wall clock speedups~\citep{meterez2025seesaw}.
%
%We also plan to explore dynamic measurement of the gradient noise
%scale~\citep{mccandlish2018empirical}, which has been shown to be
%correlated with the empirical $\cbs$, and may thus help guide a
%dynamic batch size schedule.

%- We also only consider training TIME and training COMPUTE, rather
%than inference TIME and inference COMPUTE.
%
%- overtrained models do take more training compute but reduce
%inference time and compute~\citep{touvron2023llama}, but we note we
%provide nuance to this idea by noting smaller can also reduce training
%TIME both because they use less time per optimization steps, and
%because they allow greater parallelism through a larger batch size.
%
%- Ideally, if we had an idea of our total expected training compute +
%inference compute needs, we could add inference compute to the y-axis
%on our Pareto plots and compare against training time, perhaps also
%incorporating any constraints we have on inference time.
